\section{Introduction}

% Opening 
Modelling heterogeneous treatment effects is challenging, especially with survival outcomes and high-dimensional covariates that may confound or modify treatment effects.
We propose a Bayesian regression tree ensemble approach that uses a horseshoe prior to shrink tree step heights.
This framework flexibly captures complex non-linear effects and interactions while providing adaptive shrinkage suited to high-dimensional causal inference with survival outcomes.

% Causal inference in high-dimensions

The conditional average treatment effect (CATE) quantifies how treatment effects vary across covariate profiles and is a central object in modelling heterogeneous treatment effects \citep{Caron2022CATE}.
Estimating the CATE requires careful adjustment for confounding variables.
This task becomes particularly challenging in high-dimensional settings with many covariates and complex interactions. 
In contrast, the average treatment effect (ATE) is low-dimensional and has been widely studied in high-dimensional contexts \citep{Farrell2015RobustATE, Antonelli2018DRMatching, Antonelli2020Averaging, Ning2020HDPS, Antonelli2022HighDimCausal}.
Researchers can treat high-dimensional components as nuisance terms when estimating the ATE, which simplifies inference.
CATE estimation is fundamentally more difficult because the parameter of interest is high-dimensional and depends directly on covariates \citep{Alaa2018LimitsHTE}.
Recent work has proposed flexible and robust methods for estimating heterogeneous treatment effects in high-dimensional covariate spaces \citep{Powers2018HTE, Fan2020CATE, Shin2023DRInference}.
Despite this progress, a gap remains in extending these high-dimensional causal inference methods to survival data, where censoring and time-to-event outcomes introduce additional complexity.

% Tree-based method for causal inference
Tree-based methods are widely used for causal inference in both frequentist and Bayesian frameworks. 
On the frequentist side, the causal forest algorithm of \citet{Wager2018CausalForest} extends the sample-splitting approach of \citet{Athey2016CausalTrees} to estimate heterogeneous treatment effects non-parametrically. 
This approach has been adapted to survival data by \citet{Cui2023CausalSurvivalForest}. 
In the Bayesian setting, the Bayesian causal forest approach of \citet{hahn2020bayesian} extends earlier work by \citet{Hill2011}. 
Hill applied Bayesian additive regression trees (BART) \citep{BARToriginal} to estimate treatment effect heterogeneity.
BART-based methods have demonstrated strong empirical performance and often outperform alternative estimators in benchmarking studies \citep{Dorie2019CausalCompetition, Thal2023ACIC}. 
Researchers have adapted these models to survival data to estimate conditional average treatment effects \citep{Bonato2011BayesianEnsemble, henderson2020individualized, Hu2023FlexibleCausal, Sun2025AFTCure}.
Recent work shows that AFT-BART \citep{henderson2020individualized} outperforms other causal machine learning approaches in survival settings \citep{Hu2021HeterogeneousSurvival}.

% Adaptation of BART to high-dimensional settings
We adapt Bayesian additive regression trees to high-dimensional data by placing a global--local shrinkage prior on the step heights.
In the original BART formulation, regularisation is imposed primarily through constraints on the tree structure, such as limiting tree depth or controlling the likelihood of node splits \citep{BARToriginal}.
More recent extensions, such as the Dirichlet additive regression trees (DART) of \citet{linero2018dirichlet}, introduced Dirichlet priors on splitting proportions to encourage sparsity in the covariate space.
\citet{caron2022shrinkage} further extended this idea to causal inference. 
Our approach departs from this line of work by shifting the focus of regularisation away from the tree structure and onto the leaf node parameters.
This approach preserves the flexibility of BART to model complex non-linear relationships. 
At the same time, it introduces adaptive shrinkage that suits high-dimensional causal inference. 
Rather than enforcing exact sparsity, our method applies continuous shrinkage that down-weights weak effects while still allowing strong signals to be detected when supported by the data.


% Importance of unconfoundedness in causal modelling
Causal inference typically relies on the \textit{unconfoundedness} assumption to identify treatment effects from observational data \citep{ImbensRubin2015}.
It requires that, conditional on observed covariates $X$, treatment assignment $A$ is independent of the potential outcomes $T(a)$:
\begin{equation}\label{unconfoundedness}
T(a) \ind A \mid X = x \quad \text{for } a \in \{0, 1\}.
\end{equation}
Valid causal inference is possible by adjusting for the relevant covariates if this condition holds.
Violating unconfoundedness may lead to biased estimates \citep{Chernozhukov2022OVBCML}.

In high-dimensional settings such adjustment is inherently affected by regularisation, as the estimand itself is a function of $X$.
Although treatment effect heterogeneity may be complex, it is typically assumed to be driven by an unknown low-dimensional structure in the covariates, such as a small subset of effect modifiers and confounders.
When $p$ is large relative to $n$, effective adjustment may benefit from regularisation to control model
complexity and to distinguish signal from noise.



\cite{hahn2018regularization} show that regularisation can induce bias even when the outcome model is correctly specified and all confounders are observed. 
This phenomenon, termed \textit{regularisation-induced confounding} (RIC), arises when shrinkage attenuates adjustment for covariates associated with treatment assignment, thereby distorting the estimated treatment effect. 
The risk in high-dimensional settings is particularly pronounced for weak confounders that strongly predict treatment but are only weakly associated with the outcome.
We adopt a modelling strategy that retains all covariates and controls complexity through adaptive shrinkage rather than variable exclusion.
We investigate the resulting bias empirically in simulation settings designed to exhibit RIC \citep{hahn2020bayesian}.
Motivated by this interaction between regularisation and confounding, we employ global--local shrinkage priors for adaptive regularisation of treatment effect heterogeneity.


% Bayesian shrinkage priors
% Bayesian shrinkage priors
Global--local shrinkage priors provide a fully Bayesian and computationally attractive alternative to exact sparsity-inducing methods \citep{Mitchell1988SpikeSlab, Liu2021ABC}. 
Originally developed for sparse regression problems \citep{OriginalHorseshoe}, these priors are now used across a wide range of contexts \citep{Kowal2019DynamicShrinkage, Li2019GraphicalHorseshoe, Louizos2017BayesianCompression} and admit many formulations \citep{tipping2001, johnstone2004, Bhadra2019}\citep{park2008}. 
Among these, the horseshoe prior has emerged as a particularly attractive choice because of its strong theoretical properties and empirical performance in high-dimensional settings \citep{vdPas2017Horseshoe, hahn2020bayesian, DattaGhosh2013, vdPas2014, vdPas2017, bhadra2019lasso}. 
We focus on the horseshoe prior to introduce continuous shrinkage directly on the tree step heights.
This approach retains all covariates while adaptively shrinking their contributions. 
The proposed methodology is readily extendable to other global--local shrinkage priors that allow Gibbs sampling and could be further generalised to incorporate different shrinkage strategies with additional computational effort.

% Our ocntributions
Our work contributes to the literature in several ways:
\begin{enumerate}
    \item We introduce a novel regularisation strategy for Bayesian Additive Regression Trees based on the step-height prior, where we use the horseshoe prior to induce continuous shrinkage rather than relying solely on regularisation via the tree structure.
    \item We develop a computational framework for fitting these models via reversible jump Metropolis-Hastings. The implementation is publicly available in the \texttt{R} package \texttt{ShrinkageTrees} \citep{ShrinkageTreesR}.
    \item We demonstrate how this framework can be effectively applied to estimate heterogeneous treatment effects in survival data, with a particular focus on high-dimensional covariates.
\end{enumerate}

% Organisation of the paper
The remainder of this paper is organised as follows.
In Section~\ref{sec:context} we describe the causal framework and notation.
In Section~\ref{section:Model} we introduce the model formulation and prior specification. 
In Section~\ref{sec:sampling} we describe the computational strategy and posterior inference procedure. 
We present the simulation results in Section~\ref{section:Simulations}.
Section~\ref{section:PDAC} provides a detailed case study of pancreatic cancer data. 
We conclude with a discussion in Section~\ref{section:Discussion}.



